\subsubsection*{3.1.11.}
实际中，常将 $T$ 的计算分为两步。

\begin{enumerate}[label=(\alph*)]
\item 固定满足 $t\leq\tau(d)$ 的 $t$ 和 $d$，选择 $y\in D_t^*$，并在 $f^{-1}(S_{|y|})\cap B_{t,d}$ 上选择一个联络 $\nabla$。此联络定义
\[
  V_0=f^{-1}(y)\cap B_{t,d}
  \]
到自身的微分同胚 $T_0$，但它未必在边界上为恒等。
\item 选择 $D_t$ 上 $\partial B_{t,d}$ 的一个平凡化。积分 $\nabla$ 给出从 $\partial V$ 到
\[
  f^{-1}(e^{2\pi i\theta}y)\cap\partial B_{t,d},
  \]
的微分同胚 $\varphi(\theta)$，目标通过平凡化与 $\partial V$ 等同。对 $0\leq\theta\leq1$，$\varphi(\theta)$ 是 $\operatorname{Diff}(\partial V)$ 中从恒等映射连至 $\partial T_0$ 的道路。
\end{enumerate}

所得有序对 $(T_0,\varphi)$ 依赖于所作的任意选择。由于这些选择构成的空间可缩，它在同伦意义下不依赖于这些选择。对 3.1.3 中考虑的变换 $T$，$(T,\text{常值道路})$ 是这里所考虑的一对。

\subsubsection*{3.1.12.}
映射
\[
\operatorname{Diff}(V)\longrightarrow\operatorname{Diff}(\partial V)
\]
是纤维化。因此，赋值 $T\mapsto(T,\text{常值道路})$ 诱导以下两者间的双射：

\begin{enumerate}[label=\alph*)]
\item 满足 $\partial T=\operatorname{id}$ 的微分同胚 $T:V\to V$，相差同痕视为相同；以及
\item 有序对 $(T_0,\varphi)$，其中 $T_0:V\xrightarrow{\sim}V$，$\varphi$ 是从 $\operatorname{id}$ 到 $\partial T_0$ 的道路的同伦类，相差同痕视为相同。
\end{enumerate}

因此，由 3.1.11(a),(b) 算出的一对 $(T_0,\varphi)$ 确定 $T$。具体地，由 $(T_0,\varphi)$ 计算 $T$ 如下。

\medskip
\noindent\textup{($\alpha$)} 选择领圈嵌入
\[
  [0,1]\times\partial V\xhookrightarrow{\ \alpha\ }V,
  \qquad \alpha(1,x)=x,
  \]
及一个代表
$\varphi:[0,1]\to\operatorname{Diff}(\partial V)$
使将其在负参数处定义为 $\operatorname{id}$ 所得的延拓为 $C^\infty$ 的。在领圈上令
\[
  \beta(x)=\operatorname{pr}_1\alpha^{-1}(x).
  \]
于是 $\beta$ 无临界点，取值于 $[0,1]$，且在 $\partial V$ 上等于 $1$。

\medskip
\noindent\textup{($\beta$)} 在领圈之外取 $T=T_0$，在领圈上令
\[
  T(x)=T_0\bigl(\varphi(\beta(x))^{-1}(x)\bigr).
  \]

\subsubsection*{3.1.13.}
下面是 3.1.11(b) 的一种变体，用以计算道路 $\varphi$ 的同伦类。在每条射线 $\arg(z)=\theta$ 上，选择丛 $\partial B_{t,d}$ 的一个随 $\theta$ 连续变化的平凡化。使用过 $y$ 的射线上的平凡化，将
\[
(\partial V)_0=f^{-1}(0)\cap\partial B_{t,d}
\]
等同于 $\partial V$。设
\[
\Psi_\theta:V\longrightarrow
f^{-1}(e^{i\theta}y)\cap B_{t,d}
\]
由积分 $\nabla$ 得到。将 $\partial\Psi_\theta$ 沿射线搬运，得到 $(\partial V)_0$、从而 $\partial V$ 到自身的微分同胚。这就是 $\varphi(\theta)$。